\documentclass[11pt,a4paper]{article}
\usepackage[utf8]{inputenc}
\usepackage[T1]{fontenc}
\usepackage[ngerman]{babel}
\usepackage{geometry}
\geometry{margin=2.5cm}
\usepackage{booktabs}
\usepackage{enumitem}
\usepackage{xcolor}
\usepackage{scrextend}
\usepackage{hyperref}
\usepackage{tabularx}
\usepackage{longtable}
\hypersetup{
	colorlinks=true,
	linkcolor=blue,
	citecolor=blue,
	urlcolor=blue,
}

\title{Risikovergleich gängiger Drogen in Deutschland \\ mit Schwerpunkt Jugendschutz}
\author{Zusammenstellung auf Basis aktueller Forschung}
\date{\today}

\begin{document}
	
	\maketitle
	
	\begin{abstract}
		Die Frage nach der Gefährlichkeit von Drogen ist komplex, da verschiedene Substanzen auf unterschiedliche Weise Schaden anrichten können. Dieses Dokument vergleicht zunächst die drei in Deutschland legal verfügbaren psychoaktiven Substanzen – Tabak, Alkohol und Cannabis – und geht anschließend detailliert auf die Risiken illegaler Drogen wie Opioide, Kokain, Amphetamine und MDMA ein. Es werden wissenschaftliche Erkenntnisse zu akuter Toxizität, Langzeitschäden, Abhängigkeitspotenzial und gesellschaftlichen Folgen zusammengefasst. Ein weiterer Abschnitt beleuchtet die medizinische Anwendung dieser Substanzen. Ergänzt wird der Vergleich durch die zentralen Regelungen des deutschen Jugendschutzes sowie durch einen Überblick über Präventions- und Behandlungsmöglichkeiten.
	\end{abstract}
	
	\section{Einleitung}
	
	In der öffentlichen Diskussion werden Drogen häufig pauschal bewertet, ohne die großen Unterschiede in ihren Risikoprofilen zu berücksichtigen. Während Alkohol und Tabak als legale Genussmittel tief in der Gesellschaft verwurzelt sind, gilt Cannabis oft als relativ harmlos. Ein direkter Vergleich zeigt jedoch, dass Alkohol und Tabak in vielen Kategorien ein höheres Schadenspotenzial aufweisen als Cannabis. Darüber hinaus existieren illegale Substanzen mit teils extrem hohen Risiken, insbesondere für akute Überdosierungen und schwere Abhängigkeiten. Ziel dieses Dokuments ist es, eine differenzierte, wissenschaftlich fundierte Übersicht über die Risiken der verbreitetsten Drogen in Deutschland zu geben – geordnet nach ihrem Rechtsstatus. Ein besonderer Fokus liegt auf dem Schutz von Kindern und Jugendlichen durch gesetzliche Regelungen sowie auf Präventions- und Behandlungsangeboten. Zudem wird das therapeutische Potenzial vieler dieser Substanzen dargestellt, das in kontrollierten medizinischen Kontexten genutzt wird.
	
	\section{Legale Drogen: Tabak, Alkohol und Cannabis}
	
	Diese drei Substanzen sind in Deutschland für Erwachsene legal erhältlich (Cannabis seit 2024 unter bestimmten Bedingungen). Sie unterscheiden sich jedoch erheblich in ihrem Gefährdungspotenzial.
	
	\subsection{Tabak / Nikotin}
	
	Tabak ist die mit Abstand tödlichste legale Droge in Deutschland. \cite{dhs-tabak}
	
	\begin{itemize}[leftmargin=*, noitemsep]
		\item \textbf{Akute Toxizität:} Gering – eine akute Nikotinvergiftung ist durch normales Rauchen kaum möglich.
		\item \textbf{Organschäden:} Schwerwiegend – chronisch obstruktive Lungenerkrankung (COPD), Lungenemphysem, Herz-Kreislauf-Erkrankungen (Herzinfarkt, Schlaganfall), periphere arterielle Verschlusskrankheit.
		\item \textbf{Krebsrisiko:} Extrem hoch – Tabak ist ein Gruppe-1-Karzinogen und verursacht Lungen-, Mund-, Rachen-, Speiseröhren-, Blasen-, Nieren- und Bauchspeicheldrüsenkrebs.
		\item \textbf{Psychische Risiken:} Hohe Abhängigkeit (körperlich und psychisch); Entzugserscheinungen (Reizbarkeit, Schlafstörungen, Konzentrationsprobleme). Keine direkte Psychosegefahr.
		\item \textbf{Abhängigkeit:} Extrem hoch – etwa 80\% der Raucher sind nikotinabhängig. Die Bindung an die Substanz ist mit der von Heroin vergleichbar.
		\item \textbf{Gesellschaftlicher Schaden:} Sehr hoch – jährlich über 127.000 Todesfälle in Deutschland, volkswirtschaftliche Kosten von ca. 80 Milliarden Euro pro Jahr (durch Produktivitätsverluste, Krankenkosten, Frühberentung).
	\end{itemize}
	
	\subsection{Alkohol}
	
	Alkohol ist ein Zellgift und in nahezu jeder Risikokategorie gefährlicher als Cannabis. \cite{who-carcinogen}
	
	\begin{itemize}[leftmargin=*, noitemsep]
		\item \textbf{Akute Toxizität:} Hoch – eine Alkoholvergiftung kann tödlich enden. In Deutschland gibt es jährlich ca. 44.000 alkoholbedingte Todesfälle. \cite{who-deaths}\cite{dhs-deaths}
		\item \textbf{Organschäden:} Schwer – Leberzirrhose, Pankreatitis, alkoholische Kardiomyopathie, Hirnatrophie.
		\item \textbf{Krebsrisiko:} Eindeutig (Gruppe 1) – Alkohol ist kausal mit mindestens sieben Krebsarten verbunden (Mund, Rachen, Speiseröhre, Leber, Kehlkopf, Darm, Brust). \cite{who-carcinogen}
		\item \textbf{Psychische Risiken:} Alkoholbedingte Demenz, Depressionen, Angststörungen. Körperlicher Entzug kann lebensbedrohlich sein (Delirium tremens).
		\item \textbf{Abhängigkeit:} Hoch (körperlich und psychisch) – etwa 1,6 Millionen Menschen in Deutschland sind alkoholabhängig. \cite{dhs-abhaengigkeit}
		\item \textbf{Gesellschaftlicher Schaden:} Sehr hoch – 57 Milliarden Euro volkswirtschaftliche Kosten pro Jahr, 30\% aller tödlichen Verkehrsunfälle, hohe Gewalt- und Kriminalitätsbelastung. \cite{dhs-costs}
	\end{itemize}
	
	\subsection{Cannabis}
	
	Cannabis weist im Vergleich zu Alkohol und Tabak ein deutlich geringeres Risiko für körperliche Schäden und gesellschaftliche Folgen auf, birgt jedoch spezifische psychische Gefahren. \cite{nida-no-overdose}
	
	\begin{itemize}[leftmargin=*, noitemsep]
		\item \textbf{Akute Toxizität:} Sehr gering – eine tödliche Überdosis ist extrem unwahrscheinlich. In Deutschland sind keine entsprechenden Fälle verzeichnet. \cite{bmg-cannabis}
		\item \textbf{Organschäden:} Gering – Hauptrisiko durch das Rauchen (oft mit Tabak gemischt) führt zu Atemwegsreizungen. Hinweise auf ein erhöhtes Herz-Kreislauf-Risiko, aber kein eindeutiger Nachweis für Krebs. \cite{nida-heart}
		\item \textbf{Krebsrisiko:} Kein eindeutiger Nachweis für Cannabis allein; Risiko durch Tabak-Mischkonsum.
		\item \textbf{Psychische Risiken:} Erhöhtes Psychoserisiko – besonders bei Jugendlichen mit entsprechender Veranlagung (bis zu 11-fach erhöht). \cite{cannabis-psychosis} Hochpotentes Cannabis verdoppelt das Risiko für psychotische Symptome zusätzlich. \cite{cannabis-high-potency} Beeinträchtigung der Gehirnentwicklung bei Jugendlichen.
		\item \textbf{Abhängigkeit:} Mittel, vorwiegend psychisch – etwa 9\% aller Konsumenten entwickeln eine Abhängigkeit, bei Jugendlichen bis zu 17\%. \cite{dhs-cannabis-abhaengigkeit}
		\item \textbf{Gesellschaftlicher Schaden:} Deutlich geringer als bei Alkohol oder Tabak.
	\end{itemize}
	
	\subsection{Vergleichstabelle: Legale Drogen}
	
	\begin{table}[h!]
		\centering
		\small
		\begin{tabularx}{\textwidth}{@{} l X X X @{}}
			\toprule
			\textbf{Kategorie} & \textbf{Tabak / Nikotin} & \textbf{Alkohol} & \textbf{Cannabis} \\ \midrule
			Akute Toxizität & Gering & Hoch (tödliche Überdosis möglich) & Sehr gering (keine tödliche Überdosis) \\
			Organschäden & Schwer (Lunge, Herz-Kreislauf) & Schwer (Leber, Herz, Gehirn, Pankreas) & Gering (Hauptrisiko durch Rauchen) \\
			Krebsrisiko & Extrem hoch (Gruppe 1) & Eindeutig krebserregend (Gruppe 1) & Kein eindeutiger Nachweis \\
			Psychische Risiken & Abhängigkeit, Entzug & Demenz, Depression, lebensgefährlicher Entzug & Erhöhtes Psychoserisiko (besonders bei Jugendlichen) \\
			Abhängigkeit & Extrem hoch (körperlich \& psychisch) & Hoch (körperlich \& psychisch) & Mittel (vorwiegend psychisch, 9-17\%) \\
			Gesellschaftsschaden & Sehr hoch (80 Mrd. €/Jahr) & Sehr hoch (57 Mrd. €/Jahr) & Gering im Vergleich \\
			\bottomrule
		\end{tabularx}
	\end{table}
	
	\section{Illegale Drogen: Detailierte Betrachtung}
	
	Die folgenden Substanzen unterliegen in Deutschland dem Betäubungsmittelgesetz (BtMG). Ihr Erwerb, Besitz und Handel sind für alle Altersgruppen strafbar. Die Risikoprofile variieren stark – von extrem suchtgefährdenden Opioiden bis zu eher moderat gefährlichen Substanzen wie MDMA.
	
	\subsection{Opioide (Heroin, Fentanyl, synthetische Opioide)}
	
	Opioide verursachen den Großteil der drogenbedingten Todesfälle in Deutschland (ca. 80\%). \cite{drugcom-opioid}
	
	\begin{itemize}[leftmargin=*, noitemsep]
		\item \textbf{Akute Toxizität:} Extrem hoch – Atemstillstand bereits bei geringer Überdosierung. Fentanyl ist etwa 50-100 Mal potenter als Heroin.
		\item \textbf{Organschäden:} Schäden durch intravenösen Konsum (Hepatitis C, HIV, Abszesse, Endokarditis). Bei sauberer Substanz und Applikation keine direkte Organtoxizität.
		\item \textbf{Krebsrisiko:} Kein direkt karzinogen.
		\item \textbf{Psychische Risiken:} Schwerste Abhängigkeit (körperlich dominant). Entzugssyndrom mit heftigen körperlichen Symptomen (Schmerzen, Durchfall, Erbrechen, Krämpfe) – jedoch nicht lebensbedrohlich.
		\item \textbf{Abhängigkeit:} Extrem hoch – bereits nach kurzem regelmäßigem Konsum tritt eine starke körperliche Abhängigkeit ein.
		\item \textbf{Gesellschaftlicher Schaden:} Hoch – Beschaffungskriminalität, Überdosis-Todesfälle, hohe Kosten für Substitutionsbehandlung und Notfallmedizin.
	\end{itemize}
	
	\subsection{Kokain und Crack}
	
	Kokain (intranasal) und seine rauchbare Form Crack unterscheiden sich in Suchtpotenzial und gesundheitlichen Folgen. \cite{nutt-scale}
	
	\begin{itemize}[leftmargin=*, noitemsep]
		\item \textbf{Akute Toxizität:} Hoch – Herzinfarkt, Schlaganfall, Krampfanfälle, Hyperthermie. Crack führt zu rascherer und intensiverer Wirkung, damit höheres Überdosisrisiko.
		\item \textbf{Organschäden:} Herz-Kreislauf-Schäden (Kardiomyopathie, Aortendissektion), Nasenseptumperforation bei nasalem Konsum.
		\item \textbf{Krebsrisiko:} Kein eindeutig karzinogen.
		\item \textbf{Psychische Risiken:} Paranoide Psychosen, Aggressivität, schwere Depressionen nach Konsum ("Cocaine crash").
		\item \textbf{Abhängigkeit:} Sehr hoch – Crack hat ein extremes psychisches Abhängigkeitspotenzial; körperliche Entzugserscheinungen sind mild, aber psychischer Craving sehr stark.
		\item \textbf{Gesellschaftlicher Schaden:} Hoch – insbesondere Crack ist mit Beschaffungskriminalität und sozialem Verfall assoziiert.
	\end{itemize}
	
	\subsection{Amphetamine (besonders Crystal Meth)}
	
	Methamphetamin ("Crystal Meth") ist in Deutschland insbesondere in grenznahen Regionen zu Tschechien verbreitet und führt zu rascher körperlichem und psychischem Verfall. \cite{dhs-crystal}
	
	\begin{itemize}[leftmargin=*, noitemsep]
		\item \textbf{Akute Toxizität:} Hoch – Hyperthermie (Körpertemperatur >40°C), Kreislaufversagen, Krampfanfälle, Rhabdomyolyse.
		\item \textbf{Organschäden:} Schwere Neurotoxizität (Schädigung von Dopamin- und Serotoninnervonen), schwere Zahnschäden ("Meth-Mund"), Kachexie (starker Gewichtsverlust).
		\item \textbf{Krebsrisiko:} Kein direkt karzinogen.
		\item \textbf{Psychische Risiken:} Amphetamin-Psychose (paranoid, halluzinatorisch), anhaltende kognitive Defizite, starke Stimmungsschwankungen.
		\item \textbf{Abhängigkeit:} Sehr hoch – rasche Toleranzentwicklung führt zu Dosissteigerung; schwerer Entzug mit Erschöpfung, Depression, starkem Craving.
		\item \textbf{Gesellschaftlicher Schaden:} Regional sehr hoch (Sachsen, Bayern, Berlin) – Beschaffungskriminalität, Verwahrlosung, Überlastung des Gesundheits- und Justizsystems.
	\end{itemize}
	
	\subsection{MDMA (Ecstasy)}
	
	MDMA hat ein im Vergleich zu anderen illegalen Drogen geringeres Sucht- und Überdosisrisiko, dennoch bestehen erhebliche Gefahren durch Verunreinigungen und falsches Konsumverhalten. \cite{drugcom-mdma}
	
	\begin{itemize}[leftmargin=*, noitemsep]
		\item \textbf{Akute Toxizität:} Mittel – Hyperthermie, Hyponatriämie (durch zu viel Wasser ohne Elektrolyte), Serotoninsyndrom. Todesfälle sind selten, kommen aber vor (v.a. durch Überhitzung auf Partys).
		\item \textbf{Organschäden:} Mögliche Neurotoxizität bei hochdosiertem oder häufigem Konsum (Langzeitschäden an Serotoninsystemen). Leberbelastung.
		\item \textbf{Krebsrisiko:} Kein Nachweis.
		\item \textbf{Psychische Risiken:} Angst- und Panikattacken, Depersonalisationserscheinungen, "Midweek-Blues" (serotonerger Entzug mit depressiver Verstimmung).
		\item \textbf{Abhängigkeit:} Gering – keine körperliche Abhängigkeit; psychische Abhängigkeit ist selten und meist mit partyspezifischen Kontexten verbunden.
		\item \textbf{Gesellschaftlicher Schaden:} Gering im Vergleich zu Opioiden oder Kokain – Gefahren durch verunreinigte Pillen (Streckung mit PMMA, Amphetaminen, synthetischen Cathinonen).
	\end{itemize}
	
	\subsection{Vergleichstabelle: Illegale Drogen}
	
	\begin{table}[h!]
		\centering
		\small
		\renewcommand{\arraystretch}{1.2}
		\begin{tabularx}{\textwidth}{@{} l X X X X X @{}}
			\toprule
			\textbf{Kategorie} & \textbf{Opioide (Heroin)} & \textbf{Kokain / Crack} & \textbf{Amphetamine (Crystal Meth)} & \textbf{MDMA (Ecstasy)} \\ \midrule
			Akute Toxizität & Extrem hoch (Atemstillstand) & Hoch (Herzinfarkt, Schlaganfall) & Hoch (Hyperthermie, Kreislaufversagen) & Mittel (Hyperthermie, Hyponatriämie) \\
			Organschäden & Schäden durch Injektion & Herz-Kreislauf, Nasenseptum & Neurotoxizität, Zahnschäden & Mögliche Neurotoxizität, Leber \\
			Krebsrisiko & Kein direkt karzinogen & Kein eindeutig karzinogen & Kein direkt karzinogen & Kein Nachweis \\
			Psychische Risiken & Schwerste Abhängigkeit & Psychose, Aggressivität, Depression & Psychose, kognitive Defizite & Angst, Panik, Midweek-Blues \\
			Abhängigkeit & Extrem hoch (körperlich dominant) & Sehr hoch, besonders Crack & Sehr hoch (rasche Toleranz) & Gering (keine körperliche) \\
			Gesellschaftsschaden & Hoch (Kriminalität, Todesfälle) & Hoch (Beschaffungskriminalität) & Regional sehr hoch & Gering im Vergleich \\
			\bottomrule
		\end{tabularx}
	\end{table}
	
	\section{Jugendschutz in Deutschland}
	
	Der Gesetzgeber hat für Kinder und Jugendliche gestufte Schutzbestimmungen erlassen, um die gesundheitlichen Risiken und die Entwicklungsschäden durch psychoaktive Substanzen zu minimieren. Die wichtigsten Regelungen im Überblick:
	
	\subsection{Alkohol (\S~9 JuSchG)}
	\begin{itemize}[leftmargin=*, noitemsep]
		\item \textbf{Unter 14 Jahren:} Abgabe und Konsum jeglicher alkoholischer Getränke verboten.
		\item \textbf{14–15 Jahre:} Erlaubt sind Bier, Wein und Sekt ausschließlich in Begleitung einer personensorgeberechtigten Person.
		\item \textbf{16–17 Jahre:} Bier, Wein und Sekt ohne Begleitung erlaubt; branntweinhaltige Getränke (Spirituosen, alkoholhaltige Mixgetränke) bleiben verboten.
		\item \textbf{Ab 18 Jahren:} Keine Beschränkungen.
	\end{itemize}
	
	\subsection{Tabak und nikotinhaltige Erzeugnisse (\S~10 JuSchG)}
	\begin{itemize}[leftmargin=*, noitemsep]
		\item \textbf{Unter 18 Jahren:} Abgabe und Konsum von Tabakwaren, E-Zigaretten, E-Shishas sowie nikotinhaltigen Liquids grundsätzlich verboten. Das Verbot gilt auch für nikotinfreie E-Zigaretten, sofern diese zum Verdampfen von Liquids bestimmt sind.
	\end{itemize}
	
	\subsection{Cannabis (Konsumcannabisgesetz – KCanG)}
	\begin{itemize}[leftmargin=*, noitemsep]
		\item \textbf{Unter 18 Jahren:} Absolutes Cannabisverbot. Erwerb, Besitz und Anbau sind für Minderjährige strafbar. Die Weitergabe an Minderjährige wird mit Freiheitsstrafe nicht unter zwei Jahren bestraft (\S~34 KCanG).
		\item Der Konsum in unmittelbarer Gegenwart von Minderjährigen sowie in einem Umkreis von 100 Metern zu Schulen, Kindertagesstätten und Spielplätzen ist für Erwachsene verboten (\S~5 KCanG).
		\item Verstöße von Jugendlichen werden an das Jugendamt gemeldet und können zur Anordnung von Frühinterventionsmaßnahmen führen.
	\end{itemize}
	
	\subsection{Illegale Drogen (Betäubungsmittelgesetz – BtMG)}
	\begin{itemize}[leftmargin=*, noitemsep]
		\item Substanzen wie Heroin, Kokain, Amphetamine, Methamphetamin und MDMA unterliegen dem BtMG. Erwerb, Besitz und Handel sind für alle Altersgruppen strafbar.
		\item Die Abgabe von Betäubungsmitteln an Personen unter 18 Jahren oder das Bestimmen eines Minderjährigen zum unerlaubten Umgang mit BtM wird als besonders schwerer Fall geahndet (\S~29a Abs.~1 Nr.~1, \S~30 Abs.~1 Nr.~2 BtMG) und mit Freiheitsstrafe nicht unter einem Jahr bzw. nicht unter zwei Jahren bestraft.
	\end{itemize}
	
	\section{Behandlung und Prävention}
	
	Neben gesetzlichen Verboten und Altersgrenzen sind gezielte Präventionsmaßnahmen und ein niedrigschwelliges Behandlungssystem zentrale Säulen der Suchtpolitik. Sie zielen darauf ab, den Einstieg in den Konsum zu verzögern bzw. zu verhindern und betroffenen Personen einen Ausweg aus der Abhängigkeit zu ermöglichen\cite{dhs-praevention}.
	
	\subsection{Präventionsstrategien}
	\begin{itemize}[leftmargin=*, noitemsep]
		\item \textbf{Primärprävention:} Universelle Aufklärung in Schulen, Kindertagesstätten und der Jugendhilfe (z.B. das Programm „Klasse2000“ oder „HaLT – Hart am Limit“). Dazu gehören auch strukturelle Maßnahmen wie Werbebeschränkungen, Steuerpolitik und die Begrenzung von Verkaufsstellen\cite{dhs-praevention}.
		\item \textbf{Selektive Prävention:} Angebote für Risikogruppen, etwa Kinder suchtkranker Eltern, Jugendliche mit Verhaltensauffälligkeiten oder Auszubildende in risikobehafteten Branchen (Gastronomie, Baugewerbe).
		\item \textbf{Indizierte Prävention:} Frühintervention bei erstem problematischem Konsum, z.B. durch Kurzberatung in der Hausarztpraxis, im Jugendamt oder über die Drogenhilfe. Hierzu zählen auch jugendgerichtsnahe Maßnahmen wie die Teilnahme an suchtpräventiven Kursen nach \S~10 JGG.
	\end{itemize}
	
	\subsection{Behandlungsmöglichkeiten}
	\begin{itemize}[leftmargin=*, noitemsep]
		\item \textbf{Entgiftung (qualifizierter Entzug):} Medizinisch überwachte Ausleitung bei körperlicher Abhängigkeit (Alkohol, Opioide, Benzodiazepine). Notwendig wegen möglicher Komplikationen wie Delirium tremens oder Krampfanfällen. Stationäre Entgiftung in Krankenhäusern oder spezialisierten Abteilungen.
		\item \textbf{Psychosoziale Therapien:} Verhaltenstherapie, motivierende Gesprächsführung, kontingenzmanagementbasierte Ansätze. Sowohl ambulant als auch stationär verfügbar.
		\item \textbf{Medikamentengestützte Behandlung:} 
		\begin{itemize}
			\item \textit{Opioide:} Substitutionsbehandlung mit Methadon, Buprenorphin oder retardiertem Morphin (substitutionsgestützte Behandlung).
			\item \textit{Alkohol:} Rückfallprophylaktika (Naltrexon, Nalmefen, Acamprosat) sowie aversive Substanzen (Disulfiram) unter Aufsicht.
			\item \textit{Tabak:} Nikotinersatztherapie (Pflaster, Kaugummi), Vareniclin, Bupropion.
		\end{itemize}
		\item \textbf{Rehabilitation:} Stationäre oder ambulante Langzeittherapie in spezialisierten Einrichtungen (Fachkliniken für Sucht). Ziel ist die dauerhafte Abstinenz oder die Reduktion von Konsumfolgen.
		\item \textbf{Selbsthilfe und Beratung:} Anonyme Alkoholiker, Kreuzbund, Guttempler, cannabisbezogene Selbsthilfegruppen sowie Suchtberatungsstellen der Caritas, Diakonie und kommunaler Träger.
	\end{itemize}
	
	Eine frühzeitige Intervention verbessert die Prognose erheblich. Die gesetzlichen Krankenkassen übernehmen bei gesicherter Diagnose die Kosten für Entgiftung und Rehabilitation\cite{dhs-behandlung}.
	
	\section{Medizinische Anwendung von Drogen}
	
	Die bisherige Betrachtung fokussierte auf die Risiken psychoaktiver Substanzen. Viele dieser Stoffe besitzen jedoch ein erhebliches therapeutisches Potenzial und werden seit Jahrzehnten oder werden zunehmend in kontrollierten medizinischen Kontexten eingesetzt. Der folgende Abschnitt beleuchtet den Stand der medizinischen Anwendung der zuvor besprochenen Substanzen.
	
	\subsection{Cannabis: Etablierte Medizin bei strengen Regeln}
	
	Medizinisches Cannabis hat inzwischen einen festen Platz in der deutschen Therapielandschaft, insbesondere bei chronischen Schmerzen, Spastiken (z.B. bei Multipler Sklerose), Übelkeit und Appetitlosigkeit (z.B. unter Chemotherapie). Seit 2017 können Ärzte es auf Kassenrezept verschreiben, wenn eine schwerwiegende Erkrankung vorliegt, keine Alternative verfügbar ist und eine Verbesserung zu erwarten ist. Aktuelle Daten zeigen, dass über drei Viertel der Verordnungen auf chronische Schmerzen entfallen. Aufgrund eines starken Anstiegs der Verordnungen, besonders über telemedizinische Plattformen, wurden die Regeln 2025 verschärft: Eine Erstverschreibung ist nun nur noch nach persönlichem Arztkontakt möglich, und der Versandhandel mit Cannabisblüten ist verboten.\cite{medical-cannabis}
	
	\subsection{Ketamin \& Esketamin: Durchbruch in der Depressionsbehandlung}
	
	Ein Meilenstein der letzten Jahre ist die Zulassung von Esketamin (einem Ketamin-Abkömmling) als Nasenspray zur Behandlung schwerer, therapieresistenter Depressionen. Dieser Wirkstoff ist ein Paradebeispiel für den Wandel einer Partydroge zum Therapeutikum. Im Gegensatz zu herkömmlichen Antidepressiva, die oft Wochen brauchen, kann die antidepressive Wirkung von Esketamin innerhalb weniger Stunden eintreten, was insbesondere in akuten suizidalen Krisen lebensrettend sein kann. Die Behandlung findet unter strenger ärztlicher Aufsicht statt und ist Patient:innen vorbehalten, bei denen andere Therapien versagt haben.\cite{esketamine}
	
	\subsection{MDMA, Psilocybin \& LSD: Die Renaissance der Psychedelika}
	
	Eine der aufregendsten Entwicklungen ist die Wiederentdeckung von Psychedelika für die Psychotherapie. \textbf{MDMA} wird in klinischen Studien mit großem Erfolg zur Behandlung von posttraumatischen Belastungsstörungen (PTBS) eingesetzt. Es wirkt angstlösend und fördert die therapeutische Beziehung, wodurch Trauma-Erinnerungen sicherer verarbeitet werden können. \cite{mdma-ptsd} Der Wirkstoff aus „Magic Mushrooms“, \textbf{Psilocybin}, fördert die Bildung neuer Nervenverbindungen (Neuroplastizität). Studien zeigen sein Potenzial bei Depressionen, Angstzuständen, Suchterkrankungen und Essstörungen. \textbf{LSD} wird ähnlich erforscht, unter anderem bei Clusterkopfschmerz und Suchterkrankungen. In Deutschland ist eine Psilocybin-Therapie seit Kurzem in Härtefällen (Compassionate Use) an spezialisierten Zentren möglich.\cite{psychedelics}
	
	\subsection{Amphetamine: Standard bei ADHS und Narkolepsie}
	
	Amphetamine haben eine etablierte Rolle in der Schulmedizin. Sie sind ein wichtiger Bestandteil in der Behandlung der Aufmerksamkeitsdefizit-Hyperaktivitätsstörung (ADHS), insbesondere bei schweren Verläufen, die auf andere Medikamente nicht ansprechen. Zudem werden sie zur Therapie der Narkolepsie (krankhafte Schlafsucht) eingesetzt. Sie wirken bei diesen Erkrankungen paradoxerweise beruhigend und strukturierend.\cite{amphetamine-adhd}
	
	\subsection{Opioide: Unverzichtbare, aber risikobehaftete Schmerzmittel}
	
	Trotz ihrer Gefahren als illegale Drogen sind Opioide aus der modernen Schmerzmedizin nicht wegzudenken. Sie sind bei starken und stärksten Schmerzen, etwa bei Tumorerkrankungen, unverzichtbar. Aufgrund ihres hohen Suchtpotenzials unterliegen ihre Verschreibung und Abgabe in Deutschland strengsten gesetzlichen Regelungen.\cite{opioid-pain}
	
	\subsection{Kokain: Historisches Relikt in der Medizin}
	
	Die medizinische Bedeutung von Kokain ist heute fast vollständig von historischem Wert. Es war das erste brauchbare Lokalanästhetikum, das in den 1880er Jahren eingeführt wurde. Es diente chemisch als Vorbild für viele moderne, sicherere Lokalanästhetika und wird heute nur noch in sehr seltenen Ausnahmefällen (z. B. bei bestimmten Nasen-Operationen) eingesetzt.\cite{cocaine-medical}
	
	\section{Fazit}
	
	Die Gegenüberstellung zeigt, dass die legalen Drogen Tabak und Alkohol in Deutschland die größten gesundheitlichen und gesellschaftlichen Schäden verursachen – weit mehr als Cannabis. Unter den illegalen Drogen stellen Opioide wegen ihrer extrem hohen akuten Toxizität und des starken körperlichen Abhängigkeitspotenzials die größte Gefahr dar, gefolgt von Crack und Methamphetamin. MDMA hat im Vergleich ein geringeres Sucht- und Überdosisrisiko, bleibt jedoch aufgrund von Verunreinigungen und risikoreichem Konsumverhalten (Hyperthermie, Hyponatriämie) problematisch. 
	
	Gleichzeitig zeigt der medizinische Blick, dass viele dieser Substanzen ein beachtliches therapeutisches Potenzial entfalten – von Cannabis als etabliertem Schmerzmittel über Ketamin als rapidem Antidepressivum bis hin zu psychedelischen Substanzen in der Psychotherapieforschung. Diese Doppelnatur – hohes Missbrauchs- und Schadenspotenzial bei gleichzeitiger medizinischer Nutzbarkeit – erfordert eine differenzierte Betrachtung und eine regulatorische Balance zwischen strikter Kontrolle und therapeutischer Verfügbarkeit.
	
	Keine der genannten Substanzen ist harmlos. Besonders für Kinder und Jugendliche bergen alle psychoaktiven Substanzen erhebliche Entwicklungsrisiken. Die deutschen Jugendschutzgesetze definieren daher gestaffelte Altersgrenzen für legale Drogen und sehen für illegale Drogen strikte Verbote mit empfindlichen Strafen bei Abgabe an Minderjährige vor. Ergänzend leisten Präventionsprogramme und ein vielfältiges Behandlungssystem einen wichtigen Beitrag zur Reduzierung substanzbezogener Probleme. Der sicherste Konsum ist in jedem Alter der Verzicht.
	
	\newpage
	\begin{thebibliography}{99}
		\bibitem{who-deaths} Weltgesundheitsorganisation (WHO), 2025: \href{https://www.who.int/europe/de/news/item/24-12-2025-1-in-3-injury-deaths-caused-by-alcohol--who-europe-reiterates-the-harms}{WHO/Europa weist erneut auf die schädlichen Folgen hin}.
		\bibitem{dhs-deaths} Deutsche Hauptstelle für Suchtfragen (DHS), \textit{Jahrbuch Sucht 2026}, Kapitel 2.1.
		\bibitem{nida-no-overdose} National Institute on Drug Abuse (NIDA); zitiert nach \href{https://www.dw.com/de/faktencheck-wie-ungef%C3%A4hrlich-ist-cannabis/a-65853237}{Deutsche Welle (DW), Faktencheck: Wie (un)gefährlich ist Cannabis?}
		\bibitem{bmg-cannabis} Bundesministerium für Gesundheit: \href{https://www.bundesgesundheitsministerium.de/themen/cannabis/faq-cannabis/ueberdosierung.html}{Cannabis: Fragen und Antworten zur Überdosierung}.
		\bibitem{who-carcinogen} Weltgesundheitsorganisation (WHO), 2021: \href{https://www.who.int/europe/de/home/20-10-2021-alcohol-is-one-of-the-biggest-risk-factors-for-breast-cancer}{Alkohol ist einer der größten Risikofaktoren für Brustkrebs}; IARC-Klassifizierung von 1988.
		\bibitem{nida-heart} National Institute on Drug Abuse (NIDA): \href{https://nida.nih.gov/publications/research-reports/marijuana/cardiovascular-health-effects}{Cardiovascular Health Effects of Marijuana}.
		\bibitem{cannabis-psychosis} André McDonald et al., 2025: Studie zum 11-fach erhöhten Psychoserisiko bei jugendlichen Cannabiskonsumenten, zitiert nach \href{https://www.drugcom.de/news/jugendliche-die-cannabis-konsumieren-haben-ein-erhoehtes-risiko-fuer-psychose/}{drugcom.de}.
		\bibitem{cannabis-high-potency} Lindsey Hines et al., 2024: Längsschnittstudie zu hochpotentem Cannabis und Psychoserisiko, zitiert nach \href{https://www.drugcom.de/news/hochpotenter-cannabis-verdoppelt-risiko-fuer-psychotische-symptome/}{drugcom.de}.
		\bibitem{dhs-abhaengigkeit} Deutsche Hauptstelle für Suchtfragen (DHS): \href{https://www.dhs.de/sucht/zahlen-und-fakten.html}{Zahlen und Fakten zu Alkoholabhängigkeit}.
		\bibitem{dhs-cannabis-abhaengigkeit} Deutsche Hauptstelle für Suchtfragen (DHS), \textit{Jahrbuch Sucht 2025}, Kapitel 2.5; Techniker Krankenkasse: \href{https://www.tk.de/techniker/gesundheit-und-medizin/erste-hilfe-und-gesundheitsvorsorge/sucht-und-abhaengigkeit/cannabis/cannabis-abhaengigkeitspotenzial-2075314}{Cannabis-Abhängigkeitspotenzial}.
		\bibitem{dhs-costs} Deutsche Hauptstelle für Suchtfragen (DHS): \href{https://www.dhs.de/sucht/zahlen-und-fakten.html}{Volkswirtschaftliche Kosten des Alkoholkonsums}.
		% Quellen für erweiterte Substanzen
		\bibitem{euda-2025} Europäische Beobachtungsstelle für Drogen und Drogensucht (EUDA), \textit{Europäischer Drogenbericht 2025}.
		\bibitem{dhs-drugprofiles} Deutsche Hauptstelle für Suchtfragen (DHS): \href{https://www.dhs.de/suechte}{Übersicht zu Stoffgruppen und Abhängigkeitsformen}.
		\bibitem{dhs-tabak} Deutsche Hauptstelle für Suchtfragen (DHS): \href{https://www.dhs.de/suechte/tabak}{Tabak – Zahlen und Fakten}.
		\bibitem{drugcom-opioid} \href{https://www.drugcom.de/drogen/opioide/}{drugcom.de: Opioide}.
		\bibitem{nutt-scale} David Nutt et al., 2010: Drug harms in the UK: a multicriteria decision analysis, \textit{The Lancet}.
		\bibitem{dhs-crystal} Deutsche Hauptstelle für Suchtfragen (DHS): \href{https://www.dhs.de/suechte/crystal-meth}{Crystal Meth – Basisinformationen}.
		\bibitem{drugcom-mdma} \href{https://www.drugcom.de/drogen/ecstasy/}{drugcom.de: Ecstasy (MDMA)}.
		% Quellen für Behandlung und Prävention
		\bibitem{dhs-praevention} Deutsche Hauptstelle für Suchtfragen (DHS): \href{https://www.dhs.de/praevention}{Prävention von Sucht}.
		\bibitem{dhs-behandlung} Deutsche Hauptstelle für Suchtfragen (DHS): \href{https://www.dhs.de/hilfe-und-therapie}{Hilfe und Therapie bei Sucht}.
		% Neue Quellen für medizinische Anwendung
		\bibitem{medical-cannabis} Bundesinstitut für Arzneimittel und Medizinprodukte (BfArM): \href{https://www.bfarm.de/DE/Bundesopiumstelle/Cannabis/Medizinalcannabis/_node.html}{Medizinalcannabis}.
		\bibitem{esketamine} Arzneimittelkommission der deutschen Ärzteschaft: \href{https://www.akdae.de/Arzneimitteltherapie/NA/A-2021-12/}{Esketamin bei therapieresistenter Depression}.
		\bibitem{mdma-ptsd} Multidisciplinary Association for Psychedelic Studies (MAPS): \href{https://maps.org/mdma/}{MDMA-assisted therapy for PTSD}.
		\bibitem{psychedelics} Deutsche Gesellschaft für Psychiatrie und Psychotherapie, Psychosomatik und Nervenheilkunde (DGPPN): Positionspapier zur Psychedelika-Forschung, 2024.
		\bibitem{amphetamine-adhd} Leitlinie der Deutschen Gesellschaft für Kinder- und Jugendpsychiatrie (DGKJP): ADHS im Kindes- und Jugendalter, 2023.
		\bibitem{opioid-pain} Leitlinie der Deutschen Schmerzgesellschaft: Langzeitanwendung von Opioiden bei nicht-tumorbedingten Schmerzen, 2022.
		\bibitem{cocaine-medical} Bundesopiumstelle: \href{https://www.bfarm.de/DE/Bundesopiumstelle/Betaeubungsmittel/BtM-Monographie/_node.html}{Monographie Kokain (medizinische Verwendung)}.
	\end{thebibliography}
\end{document}